% GENERATED FILE. Edit canonical lessons, then run npm run generate:books.
% canonical-source-hash: fnv1a64:40d85e7e74bd9176
% canonical-lessons: LA-C421-commercium, LA-C421-negotior, LA-C421-nundinae, LA-C421-portorium, LA-C421-lucror

\chapter{Trade}
\label{ch:la-trade}
\hlchaptermodality{421}

\begin{chapteropening}
\textbf{By the end of this chapter:} \emph{I can talk about trade and doing business, the market day, customs duty and making a profit.}

Complete the last lesson of chapter 421: I can talk about trade and doing business, the market day, customs duty and making a profit.
\end{chapteropening}

\section[commercium, commerciī]{commercium, commerciī — trade, commerce; dealings}

\label{lesson:LA-C421-commercium}



\begin{quote}
Before the new one: say the Latin for a buyer, then the Latin for a seller.
\end{quote}

\subsection*{You'll want to know: commercium, commerciī}
\textbf{commercium, commerciī} — \textquotedblleft{}trade, commerce; dealings\textquotedblright{}.

\begin{quote}
Inter duās urbēs magnum commercium erat.
\end{quote}

\begin{quote}
\emph{There was a lot of trade between the two cities.}
\end{quote}

\emph{commercium, commerciī} n., \emph{com-} + \emph{merx}, goods. Any dealings with people, not only buying and selling. English \emph{commerce}.

\begin{cousinweb}
\emph{commercium} is built from \emph{com-} \textquotedblleft{}together\textquotedblright{} and \emph{merx, mercis} \textquotedblleft{}goods\textquotedblright{} (EtymDB; Lewis \& Short). English \emph{commerce} comes from it.
\end{cousinweb}

\begin{grammarlens}[title={What you've built}]
One more word for telling, explaining and describing at B1.
\end{grammarlens}

\subsection*{Your turn}
\begin{itemize}
\raggedright
  \item \emph{Say it:} \emph{commercium, commerciī}
  \item \emph{Say it:} \emph{commercium, commerciī}, once more
  \item \emph{Say it:} a sentence of your own with \emph{commercium, commerciī}
  \item \emph{Recall:} say the Latin for a buyer, then the Latin for a seller, then say \emph{commercium, commerciī} again
\end{itemize}

\begin{tcolorbox}[breakable,colback=teal!4,colframe=teal!35!black,title={Before you move on}]
What does commercium, commerciī mean? (\textquotedblleft{}Trade, commerce; dealings\textquotedblright{}.) Say it once more.
\end{tcolorbox}
\section[negōtior, negōtiārī]{negōtior, negōtiārī — to do business, to trade}

\label{lesson:LA-C421-negotior}



\begin{quote}
Before the new one: say the Latin for a seller, then the Latin for trade.
\end{quote}

\subsection*{You'll want to know: negōtior, negōtiārī}
\textbf{negōtior, negōtiārī} — \textquotedblleft{}to do business, to trade\textquotedblright{}.

\begin{quote}
In prōvinciā negōtiābātur.
\end{quote}

\begin{quote}
\emph{He was doing business in the province.}
\end{quote}

Deponent: \emph{negōtior, negōtiārī, negōtiātus sum}, built on \emph{negōtium}, business. English \emph{negotiate}.

\begin{cousinweb}
\emph{negōtior} is built on \emph{negōtium} \textquotedblleft{}business\textquotedblright{} (EtymDB; Lewis \& Short), which is \emph{nec} \textquotedblleft{}not\textquotedblright{} and \emph{ōtium} \textquotedblleft{}leisure\textquotedblright{}. English \emph{negotiate} comes from it.
\end{cousinweb}

\begin{grammarlens}[title={What you've built}]
One more word for telling, explaining and describing at B1.
\end{grammarlens}

\subsection*{Your turn}
\begin{itemize}
\raggedright
  \item \emph{Say it:} \emph{negōtior, negōtiārī}
  \item \emph{Say it:} \emph{negōtior, negōtiārī}, once more
  \item \emph{Say it:} a sentence of your own with \emph{negōtior, negōtiārī}
  \item \emph{Recall:} say the Latin for a seller, then the Latin for trade, then say \emph{negōtior, negōtiārī} again
\end{itemize}

\begin{tcolorbox}[breakable,colback=teal!4,colframe=teal!35!black,title={Before you move on}]
What does negōtior, negōtiārī mean? (\textquotedblleft{}To do business, to trade\textquotedblright{}.) Say it once more.
\end{tcolorbox}
\section[nūndinae, nūndinārum]{nūndinae, nūndinārum — market day; a market}

\label{lesson:LA-C421-nundinae}



\begin{quote}
Before the new one: say the Latin for trade, then the Latin for to do business.
\end{quote}

\subsection*{You'll want to know: nūndinae, nūndinārum}
\textbf{nūndinae, nūndinārum} — \textquotedblleft{}market day; a market\textquotedblright{}.

\begin{quote}
Rūsticī nūndinīs in oppidum veniunt.
\end{quote}

\begin{quote}
\emph{Country people come into town on market days.}
\end{quote}

\emph{nūndinae, nūndinārum} f. pl., the ninth day, counted inclusively: every eighth day was market day.

\begin{cousinweb}
\emph{nūndinae} is built on \emph{nūndinus} \textquotedblleft{}of the ninth day\textquotedblright{} (EtymDB), from \emph{novem} \textquotedblleft{}nine\textquotedblright{} and the root of \emph{diēs} \textquotedblleft{}day\textquotedblright{} (Lewis \& Short).
\end{cousinweb}

\begin{grammarlens}[title={What you've built}]
One more word for telling, explaining and describing at B1.
\end{grammarlens}

\subsection*{Your turn}
\begin{itemize}
\raggedright
  \item \emph{Say it:} \emph{nūndinae, nūndinārum}
  \item \emph{Say it:} \emph{nūndinae, nūndinārum}, once more
  \item \emph{Say it:} a sentence of your own with \emph{nūndinae, nūndinārum}
  \item \emph{Recall:} say the Latin for trade, then the Latin for to do business, then say \emph{nūndinae, nūndinārum} again
\end{itemize}

\begin{tcolorbox}[breakable,colback=teal!4,colframe=teal!35!black,title={Before you move on}]
What does nūndinae, nūndinārum mean? (\textquotedblleft{}Market day; a market\textquotedblright{}.) Say it once more.
\end{tcolorbox}
\section[portōrium, portōriī]{portōrium, portōriī — harbour dues, customs duty}

\label{lesson:LA-C421-portorium}



\begin{quote}
Before the new one: say the Latin for to do business, then the Latin for market day.
\end{quote}

\subsection*{You'll want to know: portōrium, portōriī}
\textbf{portōrium, portōriī} — \textquotedblleft{}harbour dues, customs duty\textquotedblright{}.

\begin{quote}
Mercātōrēs portōrium solvunt.
\end{quote}

\begin{quote}
\emph{The merchants pay customs duty.}
\end{quote}

\emph{portōrium, portōriī} n., a toll on goods carried through a harbour, a gate or across a frontier.

\begin{cousinweb}
\emph{portōrium} comes from the root of \emph{portō} \textquotedblleft{}to carry\textquotedblright{} (Lewis \& Short), and EtymDB links it with \emph{porta} \textquotedblleft{}a gate\textquotedblright{}: a toll on what is carried through.
\end{cousinweb}

\begin{grammarlens}[title={What you've built}]
One more word for telling, explaining and describing at B1.
\end{grammarlens}

\subsection*{Your turn}
\begin{itemize}
\raggedright
  \item \emph{Say it:} \emph{portōrium, portōriī}
  \item \emph{Say it:} \emph{portōrium, portōriī}, once more
  \item \emph{Say it:} a sentence of your own with \emph{portōrium, portōriī}
  \item \emph{Recall:} say the Latin for to do business, then the Latin for market day, then say \emph{portōrium, portōriī} again
\end{itemize}

\begin{tcolorbox}[breakable,colback=teal!4,colframe=teal!35!black,title={Before you move on}]
What does portōrium, portōriī mean? (\textquotedblleft{}Harbour dues, customs duty\textquotedblright{}.) Say it once more.
\end{tcolorbox}
\section[lucror, lucrārī]{lucror, lucrārī — to gain, to make a profit}

\label{lesson:LA-C421-lucror}



\begin{quote}
Before the new one: say the Latin for market day, then the Latin for harbour dues.
\end{quote}

\subsection*{You'll want to know: lucror, lucrārī}
\textbf{lucror, lucrārī} — \textquotedblleft{}to gain, to make a profit\textquotedblright{}.

\begin{quote}
Ille mercātor multum lucrātus est.
\end{quote}

\begin{quote}
\emph{That merchant made a lot of profit.}
\end{quote}

Deponent: \emph{lucror, lucrārī, lucrātus sum}, built on \emph{lucrum}, gain. English \emph{lucrative}.

\begin{cousinweb}
\emph{lucror} is built on \emph{lucrum} \textquotedblleft{}gain, profit\textquotedblright{} (EtymDB; Lewis \& Short). English \emph{lucrative} comes from it.
\end{cousinweb}

\begin{grammarlens}[title={What you've built}]
One more word for telling, explaining and describing at B1.
\end{grammarlens}

\subsection*{Your turn}
\begin{itemize}
\raggedright
  \item \emph{Say it:} \emph{lucror, lucrārī}
  \item \emph{Say it:} \emph{lucror, lucrārī}, once more
  \item \emph{Say it:} a sentence of your own with \emph{lucror, lucrārī}
  \item \emph{Recall:} say the Latin for market day, then the Latin for harbour dues, then say \emph{lucror, lucrārī} again
\end{itemize}

\begin{tcolorbox}[breakable,colback=teal!4,colframe=teal!35!black,title={Before you move on}]
What does lucror, lucrārī mean? (\textquotedblleft{}To gain, to make a profit\textquotedblright{}.) Say it once more.
\end{tcolorbox}
